\documentclass[12pt,letterpaper, onecolumn]{exam}
\usepackage{amsmath}
\usepackage{amssymb}
\usepackage{amsthm}
\usepackage{graphicx}
\usepackage{mathrsfs}
\usepackage{xcolor}
\usepackage[lmargin=71pt, tmargin=1.2in]{geometry}  %For centering solution box
\usepackage[english]{babel}
\newtheorem{theorem}{Theorem}
\newtheorem{lemma}[theorem]{Lemma}
\usepackage{accents}
\usepackage{ulem} % sout() to crossing things out 
\newcommand{\ubar}[1]{\underaccent{\bar}{#1}}
\newcommand{\bs}[1]{\boldsymbol{#1}}
\newcommand{\bm}[1]{\mathbf{#1}}
\newcommand{\der}[2]{\frac{d #1}{d#2}}
\newcommand{\pder}[3]{\frac{\delta^{#3} #1}{\delta^{#3}#2}}
\newcommand{\rs}{$X_1, ..., X_n$}
\newcommand{\nsum}{\sum_{i=1}^n}
\newcommand{\nprod}{\prod_{i=1}^n}
\newcommand*{\Comb}[2]{{}^{#1}C_{#2}}%
\newcommand{\ifc}{\textrm{ if }}
\newcommand{\real}{\mathbb{R}}
\newcommand{\nat}{\mathbb{N}}
\newcommand{\power}{\mathcal{P}}
\newcommand{\rational}{\mathbb{Q}}
\newcommand{\cutelim}{\lim\limits_{n\to\infty}}
\newcommand{\salg}{$\sigma$-algebra }

\newcommand{\calA}{\mathcal{A}}
\newcommand{\calO}{\mathcal{O}}
\newcommand{\calC}{\mathcal{C}}
\newcommand{\calB}{\mathcal{B}}
\newcommand{\calI}{\mathcal{I}}
\newcommand{\calF}{\mathcal{F}}
\newcommand{\calG}{\mathcal{G}}
\newcommand{\gensig}[1]{\sigma\langle #1 \rangle} % notation for a sigma algebra generated by some set 
\newcommand{\om}{\mu^*} % outer measure 
\newcommand{\rat}{\mathbb{Q}}
\newcommand{\tribrac}[1]{\langle #1 \rangle}
\newcommand{\borel}{\mathcal{B}(\mathbb{R})}
\newcommand{\boreltwo}{\mathcal{B}(\mathbb{R}^2)}
\newcommand{\borelk}[1]{\mathcal{B}(\mathbb{R}^#1)}
\newcommand{\probspace}{(\Omega, \calF, P)}

\newcommand{\Iunion}{\bigcup_{\alpha \in I}}
\newcommand{\Iintersect}{\bigcap_{\alpha \in I}}

\newcommand{\dm}{d \mu}

\newcommand{\cutesum}[1]{\sum_{#1 = 1}^\infty}
\newcommand{\ksum}{\sum_{i = 1}^{k}}


% \chead{\hline} % Un-comment to draw line below header
\thispagestyle{empty}   %For removing header/footer from page 1

\begin{document}

\begingroup  
    \centering
    \LARGE STAT 6410\\
    \LARGE Homework 7\\[0.5em]
    \large Sarah Kilgard\par
\endgroup
\rule{\textwidth}{0.4pt}
\pointsdroppedatright   %Self-explanatory
\printanswers
\renewcommand{\solutiontitle}{\noindent\textbf{Ans:}\enspace}   %Replace "Ans:" with starting keyword in solution box

\begin{questions}
%%%%%%%%%%% Question 1 %%%%%%%%%%%%%%%%%%%%    
    \question \textbf{Chapter 3, Problem 8}

    Let $X$ be a nonnegative random variable on some probability space.
    \begin{itemize}
        \item[(a)] Show that $(EX)(E\frac{1}{X}) \geq 1$. What does this say about the correlation between $X$ and $\frac{1}{X}$.
        \item[(b)] Let $f, g: \real_{+} \to \real_{-}$ be Borel measurable and such that $f(x)g(x) \geq 1$ for all $x \in \real_{+}$. Show that $Ef(X)Eg(X) \geq 1$.
    \end{itemize}

    
    (You can assume $X > 0 \: a.e.(P)$, but it is not needed)
    \begin{solution}
    (a) \begin{proof}
        Suppose $X$ is a nonnegative random variable on some probability space.

        \textbf{Case 1:} $E(X) < \infty$\\
        \textbf{Case A:} $E(\tfrac{1}{X}) < \infty$\\
        $1 = EX\frac{1}{EX} \leq EX E(\tfrac{1}{X})$, by Jensen's Inequality since $g(x) = \frac{1}{x}$ is a convex function. 

        \textbf{Case B:} $E(\tfrac{1}{X}) = \infty$ \\
        $\infty \geq E(\tfrac{1}{X}) \geq \tfrac{1}{E(X)}$ holds. \\
        $1 = EX\frac{1}{EX} \leq EX E(\tfrac{1}{X})$
        
        \textbf{Case 2:} $E(X) = \infty$ \\
        Thus, $\frac{1}{E(X)} = 0$. Since $X \geq 0$, $E(\tfrac{1}{X}) \geq 0$. Thus, $0 = \frac{1}{E(X)} \leq E(\tfrac{1}{X})$

         $1 = EX\frac{1}{EX} \leq EX E(\tfrac{1}{X})$
    \end{proof}

    This is the formula for the correlation between $X$ and $\frac{1}{X}$: \\
    $Corr(X, \tfrac{1}{X}) = \dfrac{E(X\frac{1}{X}) - (EX) (E\tfrac{1}{X})}{\sqrt{Var(X)Var(\tfrac{1}{X})}} = \dfrac{1 - (EX) (E\tfrac{1}{X})}{\sqrt{Var(X)Var(\tfrac{1}{X})}}$

    Clearly, $\sqrt{Var(X)Var(\tfrac{1}{X})} \geq 0$, so $Corr(X, \tfrac{1}{X}) \leq 0$.
    
    (b) 
    \begin{proof}
       $X > 0$, so $f(X)g(X) \leq 1$. Thus, $g(X) \leq \frac{1}{f(X)}$. \\
       Therefore, $Ef(X)Eg(X) \leq Ef(X) E \frac{1}{f(X)} \leq 1$, from part (a).
    \end{proof}
    \end{solution}

%%%%%%%%%%% Question 2 %%%%%%%%%%%%%%%%%%%%    
    \question \textbf{Chapter 3, Problem 9}

    Prove Corollary 3.1.13 using H$\ddot{o}$lder's inequality applied to an appropriate measure space. 

    
    \begin{solution}
    \begin{proof}
Define $(\Omega, \calF, \mu)$ to be measure space such that $\Omega = \nat$, $\calF = \gensig{\{\{i\}\}_{1 \leq i \leq k}}$, and $\mu$ is the counting measure.

Let $k \in \nat$. Let $a_1, a_2, ... , a_k, b_1, b_2, ... , b_k \in \real$ and $c_1, c_2, ... , c_k \in \real_+$. 

\underline{Part (a):}

Define: $f \equiv \ksum c_i^{1/p}|a_i|I(\{i\}) = \ksum c_i^{1/p}|a_i|$ \\ and $g \equiv \ksum c_i^{1/q}|b_i|I(\{i\}) = \ksum c_i^{1/q}|b_i|$. \\ Thus, $fg = \ksum c_i|a_ib_i|I(\{i\}) = \ksum c_i|a_ib_i|$

Let $1 < p < \infty$ and $q \equiv \frac{p - 1}{p} \implies \frac{1}{q} + \frac{1}{p} = 1$
\begin{align*}
    & |fg| = \ksum c_i|a_ib_i| = \ksum c_i|a_ib_i| I(\{i\}) = \int c_i|a_ib_i| d \mu(\{i\}) \\
    & = \int |fg| d\mu(\{i\})
     \leq \left(\int |f|^p d\mu(\{i\})\right)^{1/p} \left(\int |q|^q d\mu(\{i\})\right)^{1/1} && \textrm{(By H$\ddot{o}$lder's Inequality)} \\
     & = \left(\int |a_i c_i^{1/p}|^p d\mu(\{i\})\right)^{1/p} \left(\int |b_i c_i^{1/q}|^q d\mu(\{i\})\right)^{1/q} \\
     & = \left(\ksum |a_i|^p c_i \mu(\{i\})\right)^{1/p} \left(\ksum |b_i|^q c_i \mu(\{i\})\right)^{1/q} \\
     & = \left(\ksum |a_i|^p c_i I(\{i\})\right)^{1/p} \left(\ksum |b_i|^q c_i I(\{i\})\right)^{1/q} \\ 
     & = \left(\ksum |a_i|^p c_i \mu(\{i\})\right)^{1/p} \left(\ksum |b_i|^q c_i \mu(\{i\})\right)^{1/q} \\
     & = \left(\ksum |a_i|^p c_i \right)^{1/p} \left(\ksum |b_i|^q c_i\right)^{1/q}
\end{align*}
Thus, $\ksum |a_i b_i|c_i \leq \left(\ksum |a_i|^p c_i\right)^{1/p} \left( \ksum |b_i|^q c_i \right)^{1/q}$.
\end{proof}

\underline{Part (b)}:
Define $p = 2$. Thus, $q = 2$. 

From \underline{Part (a)}, $\ksum |a_i b_i| c_i \leq \left(\ksum |a_i|^2 c_i\right)^{1/2} \left(\ksum |b_i|^2 c_i\right)^{1/2}$
    \end{solution}

%%%%%%%%%%% Question 3 %%%%%%%%%%%%%%%%%%%%    
    \question \textbf{Chapter 3, Problem 17 (a) and (b)}

    Let $X$ be a random variable on a probability space $(\Omega, \calF, \mu)$. Recall that $Eh(X) = \int h(X) d\mu$ if $h(X) \in L^1(\Omega, \calF, \mu)$.
    \begin{itemize}
        \item[(a)] Show that $(E|X|^{p_1}) \leq (E|X|^{p_2})^{p_1/p_2}$ for any $0 < p_1 < p_2 < \infty$.
        \item[(b)] Show that '=' holds in (a) iff $|X|$ is a constant a.e.($\mu$)
    \end{itemize}
    \begin{solution}
        (a)

        \begin{proof}
        Let $0 < p_1 < p_2 < \infty$ be given. 
        Define: $Z \equiv |X|^{p_1}$. Note that $|Z| = Z$. \\
        Define: $a \equiv \frac{p_2}{p_1}$ and $b \equiv \frac{a}{a - 1}$
        \begin{align*}
            & E|X|^{p_1} = E Z \leq (EZ^a)^{1/a} (E 1^b)^{1/b} && (\textrm{using H$\ddot{o}$lder's Inequality}) \\
            & = (E(|X|^{p_1})^{p_2/p_1})^{\frac{1}{p_2/p_1}} = (E|X|^{p_2})^{p_1/p_2}
        \end{align*}
        \end{proof}
        (b)
        \begin{proof}
        Let $0 < p_1 < p_2 < \infty$ be given. 
        
        From the textbook, we know that the equality holds iff, $|Z|^a = c |1|^b$ for some constant $c \in (0, \infty)$ \\
        $|Z|^a = c |1|^b \implies |X|^{p_2 / p_1} = c \implies |X| = c^{p_1 / p_2}$. $c^{p_1/p_2}$ is a constant since $p_1$ and $p_2$ are fixed values.
        \end{proof}
    \end{solution}

%%%%%%%%%%% Question 4 %%%%%%%%%%%%%%%%%%%%    
    \question \textbf{Chapter 3, Problem 18}

    Let $X$ be a nonnegative random variable
    \begin{itemize}
        \item[(a)] Show that $E(X\log(X)) \geq (EX)(E\log X).$
        \item[(b)] Show that $\sqrt{1 + (EX)^2} \leq E(\sqrt{1 + X^2}) \leq 1 + EX$. 
    \end{itemize} 
    \begin{solution}
    (a)
    \begin{proof}
        Define $g(x) = x \log(x)$. Note that $g(x)$ is convex for $x > 0$. \\ Thus,
        $E(h(X)) \leq h(EX) \: \implies \: EX \log(EX) \leq E(X \log X)$.

        Define $h(x) = \log(x)$. Note that $h(x)$ is concave for $x > 0$. \\
        Thus, $E(h(X)) \leq h(EX) \: \implies \: E(\log X) \leq \log(EX)$

        Putting this all together, we get that: 
        \[E(X\log X) \geq EX \log(EX) \geq EX E(\log X)\]
    \end{proof}
    (b) 
    \begin{proof}
    Define $g(x) = \sqrt{1 + x^2}$ and $h(x) = 1 + x$. \\
    Note that for all $x \geq 0$, $g(x) \leq h(x)$ (*). Also note that both $g(x)$ and $h(x)$ are convex functions where $x > 0$. 

    From Jensen's Inequality, we know that: $\sqrt{1 + (EX)^2} \leq E(\sqrt{1 + X^2})$. 

    From (*), we know that $E(g(X)) \leq E(h(X)) \: \implies \: E(\sqrt{1 + X^2}) \leq E(1 + X) = 1 +EX$. 

    Putting this all together: $\sqrt{1 + (EX)^2} \leq E(\sqrt{1 + X^2}) \leq 1 + EX$.
    \end{proof}
    \end{solution}

\end{questions}
\end{document}